Aquest treball final de màster aborda el desenvolupament i l'ampliació del marc de simulació dinàmica de sistemes de
potència dins de la plataforma VeraGrid, el software de codi obert d'eRoots, amb l'objectiu d'incorporar tant l'anàlisi de petit senyal com els fonaments 
de la si-\\mulació de transitoris electromagnètics (EMT). Tot el desenvolupament s'ha dut a terme en llenguatge Python, 
posant ènfasi en la modularitat i l'escalabilitat del codi amb una arquitectura clara i orientada a l'evolució de l'eina.
El treball s'emmarca en el context actual dels sistemes elèctrics, caracteritzats per una creixent presència de dispositius basats 
en electrònica de potència i per la necessitat d'eines de simulació capaces de representar fenòmens dinàmics en dife-\\rents escales temporals.

En una primera part, es desenvolupa un mòdul d'anàlisi de petita senyal per a models dinàmics en domini de valor de tensió eficaç (RMS per les seves sigles en anglès, Root Mean Square) , 
basat en la formulació simbòlica d'equacions diferencials-algebraiques i en la seva linealització al voltant del punt d'operació. 
Aquesta aproximació permet el càlcul automàtic de jacobians, autovalors i factors de participació, garantint la coherència entre la 
simulació no lineal i l'anàlisi linealitzada del sistema, i facilitant l'estudi de l'estabilitat i del comportament modal del sistema elèctric. 
Els resultats obtinguts es validen mitjançant la comparació amb ANDES, una eina de referència en l'àmbit Simbòlic-Numèric.

Posteriorment, s'estableixen els fonaments per a la simulació EMT mitjançant la implementació de models en el 
domini temporal amb variables instantànies en coordenades abc. Es desenvolupen models bàsics de generador síncron, 
línies de transmissió i càrregues, amb especial èmfasi en el model de línia basat en el model de Bergeron i 
en la teoria d'ones viatgeres. La formulació EMT s'implementa utilitzant l'algoritme de Dommel conjuntament amb el marc de treball Simbòlic-Numèric 
i esquemes d'integració trapezoïdal, assegurant una representació físicament coherent dels fenòmens electromagnètics ràpids. 
La validació es realitza mitjançant diferents casos d'estudi, que permeten analitzar la propagació d'ones i els mecanismes de 
reflexió sota diverses condicions.

Els resultats obtinguts mostren un comportament basat en la teoria clàssica dels transitoris electromagnètics i confirmen la 
correcció tant de la formulació matemàtica com de la implementació numèrica. Aquest treball estableix una base sòlida, extensible i robusta per al 
desenvolupament d'un motor de simulació EMT avançat integrat dins de VeraGrid, amb un alt potencial d'impacte tant en l'àmbit acadèmic com en aplicacions 
industrials reals.


